% REVIEW DRAFT: not an input of the current published manuscript.
% This file is not proof certification; see geometry-audit-2026-11-21.json.
\subsection{An integral replacement for the Dieudonne shortcut}
\label{subsec:sr-formal-hodge-kernel-repair}

The argument here concerns the actual formal cubical Hodge module.
It does not identify it with an integral Dieudonne lattice.
Such an identification was asserted without proof in the earlier
flat-saturation draft and is not used below.
We give the finite-kernel argument, including its identification
with the mixed term of the formal action.

We use two classical finite-group facts in their exact field
scope.  Over a perfect field of characteristic \(p>0\), a finite
group of height at most one has order \(p^{\dim\operatorname{Lie}}\);
its distribution algebra is the restricted enveloping algebra
of its Lie algebra.  For a finite commutative group \(K\), the
subalgebra of \(\mathcal O(K)\) generated by its primitive
elements is the dual distribution algebra of \(K^D[F]\).
These are Cartier's height-one theorem and its functorial
restricted-PBW construction; see
\href{https://www.jmilne.org/math/Books/iAG2022.pdf}
{Milne, \emph{Algebraic Groups}}, Propositions~11.35--11.37.
There is no exclusion of \(p=2\).
For abelian varieties we use Cartier duality for the kernel of
an isogeny, and the integral characteristic polynomial of an
endomorphism, whose constant term is its degree; see
\href{https://www.jmilne.org/math/CourseNotes/AV110.pdf}
{Milne, \emph{Abelian Varieties}}, Theorem~8.1, Theorem~9.9 and Proposition~9.20
(and Mumford, \S19, Theorem~4).
Only these field-level facts are imported.

\begin{srlemma}[A finite-group tangent inequality]
\label{lem:sr-finite-group-tangent-bound}
Let \(k\) be algebraically closed of characteristic \(p>0\),
let \(K/k\) be a finite commutative group, and let
\[
 b:\operatorname{Lie}K\times\operatorname{Lie}K^D\to k
\]
be the mixed tangent pairing of Cartier evaluation
\(K\times K^D\to\mathbb G_m\).
If \(W\) is its left radical, then
\[
       p^{\dim\operatorname{Lie}K^D+\dim W}
                              \ \mid\ \operatorname{rank}K .
 \tag{FK1}
\]
\end{srlemma}
\begin{proof}
Put \(A=\mathcal O(K)\).
A tangent vector \(z\in\operatorname{Lie}K^D\) is a character
of \(K_{k[\epsilon]/\epsilon^2}\) reducing to the identity,
hence is uniquely \(1+\epsilon f_z\).
The character equation is
\[
 \Delta f_z=f_z\otimes1+1\otimes f_z,\qquad \varepsilon(f_z)=0 .
\]
It gives a linear isomorphism
\(\operatorname{Lie}K^D\simeq\operatorname{Prim}(A)\).
A tangent vector \(v\) of \(K\) is an augmentation derivation
\(A\to k\), and \(b(v,z)=v(f_z)\).

Let \(B\subset A\) be the Hopf subalgebra generated by all
primitive elements, and let \(H=\operatorname{Spec}B\).
Here is the height-one identification used above.
The first Frobenius kernel \(K^D[F]\) has the same Lie
algebra as \(K^D\).  Its distribution algebra is generated
by that Lie algebra, by restricted PBW.  Under the injection
of distribution algebras and then into the full linear dual,
\[
 \operatorname{Dist}(K^D[F])
     \hookrightarrow\operatorname{Dist}(K^D)
     \hookrightarrow \mathcal O(K^D)^\vee=A,
\]
these generators are exactly \(\operatorname{Prim}(A)\).
The second arrow need not be surjective: distributions are
supported at the identity, whereas \(K^D\) can be disconnected.
Thus its image is \(B\), and
\(\operatorname{rank}H=p^{\dim\operatorname{Lie}K^D}\).
This also proves directly that this rank is a pure power of \(p\).

The inclusion \(B\subset A\) corresponds to a faithfully flat
finite-group quotient \(K\twoheadrightarrow H\).
Let \(J\) be its kernel.
A derivation of \(A\) vanishes on \(B\) exactly when it vanishes
on every primitive generator.  Left exactness of the Lie
functor therefore says \(\operatorname{Lie}J=W\).
The height-one group \(J[F]\) has order \(p^{\dim W}\).
The two finite quotient formulas give
\[
 \operatorname{rank}K=\operatorname{rank}H\,
                              \operatorname{rank}J,\qquad
               p^{\dim W}\mid\operatorname{rank}J.
\]
This is (FK1).
\end{proof}

\paragraph{The mixed coefficient is the Cartier pairing.}
Let \(A/k\) be an abelian variety with a rigidified line \(M\)
whose homomorphism \(\lambda_M:A\to A^\vee\) is an isomorphism.
Let an \(\mathcal O_E\)-action have Rosati involution \(a\mapsto\bar a\).
For formal coordinates at zero write \(R_a\) for the tangent
matrix, and \(C_a\) for the mixed coefficient of the normalized
cubical adjoint comparison.
The actual formal action is
\[
 \rho_H(a)(v,f)=(R_av,R_{\bar a}^{\,t}f-C_a^{\,t}v).
 \tag{FK2}
\]
We verify which finite-group pairing its coefficient computes.

Let
\(\mathcal P_M(x,y)=M_{x+y}\otimes M_x^{-1}\otimes M_y^{-1}\otimes M_0\).
The adjoint relation gives an isomorphism of biextensions
\[
               \mathcal P_M(ax,y)
                            \simeq\mathcal P_M(x,\bar a y).
 \tag{FK3}
\]
We match its normalization with the actual formal frame.
Write \(e(x,y)\) for the formal trivialization of
\(\mathcal P_M(x,y)\) used by that frame; it is rigidified by
\[
                         e(0,y)=e(x,0)=1.
\]
The actual frame supplies a global nowhere vanishing generator
\(w_a\) of the line
\[
        \mathcal P_M(x,\bar a y)\otimes
                                      \mathcal P_M(ax,y)^{-1}.
\]
Here and below the value of \(w_a\) is written in the
trivializations of the indicated line factors.
The defining correspondence in the frame, evaluated on each
finite jet, is the identity
\[
 \operatorname{ad}_a(x,y)
       =w_a(x,y)\,\frac{e(ax,y)}{e(x,\bar a y)} .
 \tag{FK3a}
\]
This is exactly the cubical comparison in the frozen source
\texttt{FBIdAd.dict\_ad}; it follows there by multiplying the
pullback of \(w_a\) by the frame of \(\mathcal P_M(ax,y)\)
and comparing with the frame of
\(\mathcal P_M(x,\bar a y)\).
Thus (FK3a) is an equality on the special fibre and its
infinitesimal charts, not an identification obtained by
inverting the residue characteristic.

For completeness, there is no free scalar in (FK3a).
Two global generators of this line differ by a unit on the
proper integral variety \(A\times A\), hence by a member of
\(k^\times\).
Putting \(x=y=0\) in (FK3a), the frame condition
\(\operatorname{ad}_a(0,0)=1\) gives \(w_a(0,0)=1\).
Restriction to either axis is then one, since it is a unit
on the proper connected variety \(A\) with value one at zero.
All biextension cocycle diagrams commute: their ratios are
units on \(A^3\), and evaluation at a zero axis fixes those
constants to one.
Consequently \(w_a\) is exactly the rigidified adjoint
biextension isomorphism.  This establishes the required
normalization in characteristic \(p\), including \(p=2\).

For \(a\ne0\), the endomorphism \(a\) is an isogeny.
Its dual kernel is identified with \(K_{\bar a}=\ker\bar a\)
by \(\lambda_M\), since
\(a^\vee\lambda_M=\lambda_M\bar a\).
Restrict the rigidified adjoint biextension isomorphism to
\(K_a\times K_{\bar a}\).
Both of its line factors have a zero argument and hence
canonical trivializations.
The resulting scalar is the character obtained by
descending the degree-zero line
\(\mathcal P_M(-,z)\) along \(a\); this is the construction
of isogeny-kernel Cartier duality.
It gives the perfect pairing
\[
                         e_a:K_a\times K_{\bar a}\to\mathbb G_m,
 \tag{FK4}
\]
possibly replaced by its inverse according to the direction
chosen for the adjoint isomorphism.
One can characterize this identification without any scalar
choice: for \(z\in K_{\bar a}\), pullback by \(a\) trivializes
\(\mathcal P_M(-,z)\), and translation by \(x\in K_a\)
acts on that trivialization through \(e_a(x,z)\).
This is the universal Cartier character.

Now (FK3a) restricts to
\(\operatorname{ad}_a(x,z)=w_a(x,z)\) on these two kernels:
both formal factors in its fraction are exactly one.
For tangent vectors \(v\in\ker R_a\) and
\(z\in\ker R_{\bar a}\), substitute
\(x=\epsilon v,\ y=\delta z\) over
\(k[\epsilon,\delta]/(\epsilon^2,\delta^2)\).
The formal axes identities remove every pure term, so
\[
 \operatorname{ad}_a(\epsilon v,\delta z)
                   =1+\epsilon\delta\,v^tC_a z.
 \tag{FK5}
\]
By the preceding restriction this is the mixed derivative
of (FK4).  Taking its inverse changes only the sign of that
derivative, and hence changes neither radical.
Thus (FK2) uses the actual Cartier mixed pairing in the
actual integral formal frame.  No comparison with an
unspecified de Rham lattice is used.

\begin{srproposition}[Proper formal kernel bound]
\label{prop:sr-proper-formal-kernel-bound}
In this situation, let \(n=[E:\mathbb Q]\), \(g=\dim A\).
For every nonzero \(a\in\mathcal O_E\),
\[
 \dim_k\ker\rho_H(a)
       \le \frac{2g}{n}\,v_p|N_{E/\mathbb Q}a|.
 \tag{FK6}
\]
If \(A\ne0\), \(2g/n\) is an integer.
\end{srproposition}
\begin{proof}
For the linear map in (FK2), projection of its kernel onto
the first coordinate has image exactly
\[
 W=\{v\in\ker R_a:\ v^tC_a z=0
                                  \text{ for every }z\in\ker R_{\bar a}\}.
\]
Indeed \(R_{\bar a}^{\,t}f=C_a^{\,t}v\) is solvable exactly
when its right side annihilates \(\ker R_{\bar a}\).
The kernel of this projection has dimension
\(\dim\ker R_{\bar a}\).
By (FK5) and Lemma~\ref{lem:sr-finite-group-tangent-bound},
\[
 \dim\ker\rho_H(a)=\dim W+\dim\operatorname{Lie}K_{\bar a}
                                   \le v_p\operatorname{rank}K_a.
 \tag{FK7}
\]

We justify the norm in the last term.
Choose an integral primitive element \(u\) of \(E/\mathbb Q\).
For any prime \(\ell\ne p\), its action on \(V_\ell A\)
satisfies its separable minimal polynomial.
The characteristic polynomial of this action lies in
\(\mathbb Z[T]\) and is independent of \(\ell\).
Its roots in an algebraic closure are therefore all conjugates of \(u\),
with a common multiplicity \(r_A\).
It follows that \(2g=n r_A\).
Since every \(a\in E\) is a rational polynomial in \(u\),
its determinant on \(V_\ell A\) is
\(N_{E/\mathbb Q}(a)^{r_A}\).
The degree-characteristic-polynomial theorem gives
\[
                         \operatorname{rank}K_a
                         =\deg(a)=|N_{E/\mathbb Q}a|^{r_A}.
\]
Together with (FK7) this proves (FK6).
When \(g=0\), the module is zero and the assertion is immediate.
\end{proof}

\paragraph{The semiabelian boundary uses the same formal action.}
We state exactly the quotient data needed.
Let \(G/k\) be semiabelian, and suppose there is an equivariant
torus extension
\(1\to T\to G\stackrel q\to A\to1\)
and an isomorphism of rigidified cubical lines
\(M_G\simeq q^*M_A\).
Require that \(\lambda_{M_A}\) be an isomorphism, that the
given action on \(A\) be the one induced by \(q\), and that
its adjoint biextension isomorphisms have Rosati involution
\(a\mapsto\bar a\).
These are the compatible quotient data furnished by the closed
polarized Mumford data, including its line and action clauses;
a bare unpolarized torus extension is insufficient.

Choose formal coordinates for \(G,A\), with tangent quotient
\(P:\operatorname{Lie}G\to\operatorname{Lie}A\).
The isomorphism \(M_G\simeq q^*M_A\) changes formal
trivializations by a unit, normalized on the zero axes.
Taking its mixed derivative \(Z\) changes the adjoint
mixed matrix by
\[
 C_{G,a}+R_{G,a}^{\,t}Z-ZR_{G,\bar a}
                                         =P^t C_{A,a}P .
 \tag{FK8}
\]
This is obtained by substituting the tangent map into the
two variables of (FK3); a mixed coefficient transforms as
\(C\mapsto P^tCP\).
The two trivialization factors contribute the two displayed
coboundary terms.
The corresponding triangular change of \((v,f)\) is invertible,
so it conjugates the actual formal actions.

In that gauge, put \(T_0=\ker P\).
The map \((v,f)\mapsto f|_{T_0}\) is equivariant from the full
formal module of \(G\) to the conjugate dual of \(T_0\).
On its kernel choose a linear section \(Q\) of \(P\).
The map
\[
                           (v,f)\longmapsto(Pv,Q^tf)
 \tag{FK9}
\]
is an equivariant surjection onto the full formal module
of \(A\), with kernel \(T_0\).
For \(f|_{T_0}=0\), the identity \(f=P^tQ^tf\) verifies
the adjoint block, and (FK8) verifies its mixed block.
This proves an actual equivariant three-step filtration:
\[
               T_0,\qquad H_A,\qquad T_0^{\vee,\bar{\ }} .
 \tag{FK10}
\]
No crystalline identification is being assumed.

The character lattice of \(T\) is a finite projective
\(\mathcal O_E\)-module, of some rank \(s\).
Thus \(\dim T=ns\).
For multiplication by \(a\ne0\) on this integral
\(\mathbb Z\)-lattice, the determinant is \(N(a)^s\).
Smith diagonalization over \(\mathbb Z_p\) gives
\[
             \dim_k\ker(a\mid T_0)
                     \le s\,v_p|N(a)| .
 \tag{FK11}
\]
The same holds on its conjugate dual, since \(N(\bar a)=N(a)\).
If the total formal dimension is \(nr\), the abelian
multiplicity in (FK6) is \(r-2s\).
On any exact sequence preserved by an endomorphism, the
dimension of its kernel is at most the sum of the
kernel dimensions on submodule and quotient.
Using (FK6), (FK10) and (FK11) proves
\[
                         \dim_k\ker\rho_{H_G}(a)
                               \le r\,v_p|N(a)|.
 \tag{FK12}
\]
This is precisely the geometric input needed in the
residue-algebra freeness criterion.

\begin{srlemma}[Norm bounds force a free residue module]
\label{lem:sr-residue-norm-freeness}
Let \(W\) be a finite-dimensional representation of
\(\mathcal O_E\) over algebraically closed \(k\) of characteristic
\(p>0\).  Suppose, for every \(a\in\mathcal O_E\),
\[
 \operatorname{char}(a\mid W)
           =\operatorname{char}(a\mid E)^r\bmod p,\qquad
             \dim\ker(a\mid W)\le r\,v_p|N(a)|\quad(a\ne0).
 \tag{FK13}
\]
Then \(W\) is free of rank \(r\) over
\(k\otimes_{\mathbb Z}\mathcal O_E\).
\end{srlemma}
\begin{proof}
Because \(k\) is algebraically closed and \(E\) is a number field,
\[
 k\otimes_{\mathbb Z}\mathcal O_E
   \simeq\prod_{\mathfrak p\mid p}
          \prod_{\sigma:\kappa(\mathfrak p)\hookrightarrow k}
                                          k[t]/(t^{e_{\mathfrak p}}).
 \tag{FK14}
\]
This follows by separating the unramified residue field and
reducing an Eisenstein uniformizer polynomial modulo \(p\).
The characteristic-polynomial hypothesis gives dimension
\(e_{\mathfrak p}r\) on each factor.
For clarity, use an element supported by the corresponding
prime in \(\mathcal O_E/p\), with residue a generator of
\(\kappa(\mathfrak p)^\times\).
The conjugates of that generator are distinct; its eigenvalue
multiplicities isolate all the residue embeddings.
This determines the dimensions as asserted.

Fix \(\mathfrak p\mid p\), with residue degree \(f_{\mathfrak p}\).
Choose \(a\in\mathcal O_E\) with valuation exactly \(1\)
at \(\mathfrak p\), and valuation \(0\) at every other
prime over \(p\), by the Chinese remainder theorem
modulo \(\mathfrak p^2\) and those other primes.
Then \(v_p|N(a)|=f_{\mathfrak p}\).
On each of the \(f_{\mathfrak p}\) factors in (FK14),
\(a\) is \(t\) times a unit.
Every Jordan block has length at most \(e_{\mathfrak p}\).
Since that factor has dimension \(e_{\mathfrak p}r\),
it contributes at least \(r\) to \(\dim\ker a\).
All other factors contribute zero.
The upper bound in (FK13) forces equality in every one
of these inequalities.  There are exactly \(r\) blocks
on each factor, and their total length forces every block
to have length \(e_{\mathfrak p}\).
Thus each factor is free of rank \(r\).
Doing this for each prime proves the lemma.
\end{proof}

For the actual formal bundle on a normal model, its characteristic
polynomial identity holds at the generic point and hence
everywhere, since the bundle is locally free over the reduced base.
At a special point with the compatible polarized quotient data
above, (FK12) and the last lemma prove its actual residue-algebra
freeness.  Lifting a free residue basis gives a map
\[
 (\mathcal O_E\otimes\mathcal O_{\rm model})^r\longrightarrow\mathcal H.
\]
Nakayama makes it surjective, and equality of ordinary ranks
makes it an isomorphism.  A geometric residue-field extension
is faithfully flat, so this conclusion descends.
At a characteristic-zero point the same conclusion follows
from the split semisimple action and the characteristic
polynomial identity.

There is also no missing global-function hypothesis in forming
these formal data on the actual model.
On every integral affine base open the generic fibre of each
power of the semiabelian family is an abelian variety.
A regular function on that power is constant at the generic
fibre.  Evaluation at the identity extends that constant to a
regular function on the base, and flatness makes the generic
fibre dense, so the two functions agree.
This proves the source's \(\mathrm{PowFns}\) property on
every base open.  It ensures that the rigidified comparison
isomorphisms used to form (FK2) satisfy their required
coherence identities.

\paragraph{Scope of this repair.}
This proves the actual formal-module kernel bound and
\(\mathcal O_E\)-freeness at the semiabelian boundary whenever
the displayed compatible polarized quotient data are present.
For an actual \(\mathrm{PELUnif}\), these are obtained by
reducing the Raynaud extension, its line isomorphism,
its \(\mathcal O_E\)-equivariance and the trace-form identity.
This is the precise source scope of
\texttt{HeightCharFreeGeometric},
\texttt{ProperMultiplicationKernel},
\texttt{AbelianQuotientBilinear}, and
\texttt{KernelRankBoundNormValuation}.
It repairs only this join in the flat-saturation draft.
The formal Frobenius exact-sequence identification,
the construction of the globally compatible full Mumford models,
and the comparison of their canonical determinant extensions
with the analytic flat frame must still be audited separately;
they are not consequences of (FK12)--(FK14).
